%% kapa: exam=1 round=25 year=2014 period=2 form=A subject=appraisal_law subject_ko=부동산관계법규 seq=3 number=3
%% source: Q-Net 감정평가사 기출문제 (cst00309) artlSeq=5101207 file=extracted/1008/5101207/2014년 제25회 감정평가사 1차시험 기출문제/2교시(A형).hwp
%% conversion: kordoc 4.24.51 -> tools/gen1.py (auto); review=proofread-sonnet; flags=none
\begin{question}{3}
국토의 계획 및 이용에 관한 법령상 용도지역별 건폐율의 상한(上限)을 비교한 것으로 옳은 것은? (단, 개별 조례의 규정은 고려하지 않음)
\begin{choices}
  \item 준공업지역 > 준주거지역 > 제3종일반주거지역 = 제2종일반주거지역
  \item 준공업지역 = 준주거지역 > 제3종일반주거지역 > 제2종일반주거지역
  \item 준공업지역 = 준주거지역 > 제2종일반주거지역 > 제3종일반주거지역
  \item 준주거지역 > 준공업지역 > 제3종일반주거지역 > 제2종일반주거지역
  \item 준주거지역 > 준공업지역 > 제3종일반주거지역 = 제2종일반주거지역
\end{choices}
\end{question}
\answer{3}
